\section{Introduction}
\label{sec:introduction}

Treedepth is the minimum height of an \emph{elimination forest}: a rooted forest on
the graph's vertices in which the endpoints of every edge are comparable in the
ancestor order. For connected graphs, every elimination forest is a tree.
Recursively querying its root and then searching each remaining component gives
a graph-search strategy. Treedepth is also characterized by vertex ranking and
symbolic elimination~\cite{BodlaenderRankings}, and it is studied directly in
parameterized algorithms~\cite{Reidl2014}. Low-treedepth decompositions are
used in the structural theory of sparse graphs. Nešetřil and Ossona de Mendez
survey these methods, and their earlier work connects them to subgraph coloring
and homomorphism bounds~\cite{LowTreeDepthSurvey,TreeDepthColoring}.

A \emph{tree decomposition} consists of a tree of vertex sets, called bags,
such that every graph edge has both endpoints in some bag and the bags
containing any fixed vertex form a connected subtree. Treewidth is the minimum
possible largest bag size minus one. Decompositions of small width support
dynamic programming for many graph problems and have uses in graph-minor
theory. For surveys of these applications, see~\cite{BodlaenderTouristGuide,
BodlaenderApplications}. A path has treewidth one, whereas its treedepth grows
with its length. We write $\td(G)$ and $\operatorname{tw}(G)$ for these
parameters.

The approximation complexity of these parameters remains poorly understood.
Bonnet, Neuen, and Soko\l owski proved that approximating
treedepth within factor $1.0003$ is NP-hard~\cite{BNS}. Bonnet proved the
corresponding $1.00005$ hardness bound for treewidth~\cite{BonnetTreewidth}.
Under the Small Set Expansion Hypothesis (SSEH), treewidth is hard to
approximate within every fixed constant factor~\cite[Theorem~1.1]{APW}. On the
algorithmic side, Anand gives an $O(\log n)$-approximation for
treedepth~\cite{Anand2026}, while Feige, Hajiaghayi, and Lee give a polynomial-time
$O(\sqrt{\log(\operatorname{tw}(G)+2)})$-approximation for
treewidth~\cite{FeigeHajiaghayiLee2005}. These results leave a substantial gap
between known algorithms and hardness. We strengthen the hardness bounds and
show that they persist on graphs of bounded maximum degree, addressing the
conjecture of Schirrmacher, Siebertz, and Vigny
~\cite[Conjecture~1]{SchirrmacherSiebertzVigny2025}.

We also study \emph{total elimination depth}, $\totd(G)$, the minimum
sum of vertex depths in an elimination forest (roots have depth one).
This is the vertex-partition-tree sum of~\cite{HogemoEtAl2021}. On a
connected graph, $\totd(G)/|V(G)|$ is the expected cost of vertex-query
search with a uniform target and unit query costs
~\cite{HogemoEtAl2021,Szyfelbein2025GraphSearch}.

\subsection{Our results and techniques}
\label{subsec:results-techniques}

We first prove the hardness gaps on unrestricted simple unweighted graphs.
Here UGC denotes the ordinary Unique Games Conjecture, SSEH the Small Set
Expansion Hypothesis, and NP-hardness means unless $\mathrm{P}=\mathrm{NP}$.

\Needspace{12\baselineskip}
\begin{theorem}[Unrestricted treedepth, treewidth, and total-depth hardness]
\label{thm:main-results}
\label{thm:total-depth-main}
There is an absolute constant $\epsilon_0>0$ such that, for each row
below and every fixed factor $r$ in the indicated range,
approximating the corresponding parameter within factor $r$ is hard
in the stated sense.
\begin{center}
\begin{tabular}{lccc}
\toprule
\textup{Hardness} & \textup{Treedepth} & \textup{Treewidth}
                  & \textup{Total elimination depth}\\
\midrule
\textup{NP-hard} & $1\leq r<18/17$ & $1\leq r<12/11$
                 & $1\leq r<1+\epsilon_0$\\
\textup{Unique Games-hard} & $1\leq r<4/3$ & $1\leq r<4/3$
                           & $1\leq r<8/7$\\
\textup{SSE-hard} & \textup{every fixed }$r\geq1$
                & \textup{every fixed }$r\geq1$
                & \textup{every fixed }$r\geq1$\\
\bottomrule
\end{tabular}
\end{center}
The reductions are deterministic.
\end{theorem}

The proofs follow three source routes. In order to prove approximation hardness of treedepth and total-depth under
UGC, we reduce from the independent-set by constructing layered clique graphs. By this we mean a graph consisting of several clique layers, connected by edges inbetween clique vertices, representing neighbouring pairs in the source graph. A large independent set in the source graph corresponds then to a shallow elimination forest, by placing its complement vertices in the upper layers and the independent set vertices in the lower layers. Conversely, if all independent sets are small, an inductive reasoning over the number of layers shows that any elimination forest must be deep. 
 
For treewidth, as a source we use inapprixambility of vertex cover on 3-partite graph, refining the construction of ~\cite[Lemma~6]{BonnetTreewidth} by replicating the copies before completing each global part to a clique, which makes
the one-copy overhead negligible. Bonnet's central-bag argument then forces
a cover from every source copy into a central bag. This very same construction will also yield hardness of approximation based on $P\neq NP$, although using a different source promise.

Under SSEH, we transfer the Raghavendra-Steurer-Tulsiani partition gap~\cite{RST}
to treedepth and total-depth, by constructing a line graph of the source graph.
In the YES case, we can then construct a shallow elimination forest corresponding to the partition.
In the NO case, every balanced separator must be large, which then becomes the lower bound on the cost of any elimination tree.
 The same separator bounds give the
total-depth gap. For treewidth, the standard such a result is already given in
~\cite[Lemma~5.3]{APW}.

Under $\mathrm{P}\ne\mathrm{NP}$, we refine the truth-block and anchor
construction of Bonnet, Neuen, and
Soko\l owski~\cite{BNS}, using the parity gap of Hazan, Safra, and
Schwartz~\cite[Theorem~10]{HazanSafraSchwartz2003}. We reduce from the parity problem. The construction encodes the variables using bicliques, so that any indepdent set can choose vertices of only one side. Moreover, we add 4 vertex per each constraint, where each vertex encodes a viable assignment consistent with the constraint. This ensures, that any indepedent set can pick at most 3 of them. Thus gives an inapproximability of indepdent set on such instances, which then transfers to treedepth, and treewidth. A separate reduction from
cubic Independent Set~\cite{AlimontiKann2000} gives the total-depth gap:
source parts become
cliques, and attached leaves make the depth depend on the deepest vertex in
each part. A matching makes the split separating one part from the other
two cheapest, so the depth sum reflects the source independent-set size.

En route, we observe that all of the previously mentioned consructions, only use clique and bi-clique gadgets to enforce certain useful properties. These gadgets, force the graph to be very dense, and in particular the degree of the produced instances is unbounded. In Section 
Section~\ref{sec:bounded-degree-hardness} our next goal becomes, to turn those instances into sparse graphs, with bounded degree. Our strategy is to replace each such dense gadget, with a bounded degree expanders. In order for such a construction to work, we need to ensure, that each such gadget contains no large \emph{holes}, which means that for all pairs of sufficiently large disjoint subsets of vertices in such a gadget are connected by at least one edge. To obtain such a family of graphs, we start with an expander, and power it sufficiently many times. Each such operation decreases the guarantee on the size of the largest possible hole.
Since such replacements only delete edges, the YES-case decompositions remain
feasible. It can be shown, that in the NO case, this incurrs only a small amount of distortion in the value of the optimum value of the objective. As a corollay we obtain that all our inapproximability results can be transferred to bounded degree instances. This resolves the treedepth conjecture
of Schirrmacher, Siebertz, and Vigny~\cite[Conjecture~1]{SchirrmacherSiebertzVigny2025}.

\subsection{Applications}

Our main application is to \emph{Trivially Perfect Completion},
where edges are added until the graph is a forest closure. Two
distinct vertices are adjacent if and only if one is an ancestor of
the other in a rooted forest. The objective is the number of added
edges. For every fixed $C\ge1$, approximating this objective within
factor $C$ is SSE-hard on graphs of maximum degree bounded by a
constant depending only on $C$
(Theorem~\ref{thm:tp-completion-sse}). The SSE proof uses the same
line graph as our SSE treedepth reduction, but needs a bound on the \emph{sum} of
elimination depths. A balanced separator, whose deletion leaves every
component with at most half the vertices, supplies this estimate,
and an adjusted-cost argument handles the offset between total depth
and the number of added edges. The identity relating these objectives
is established in~\cite{HogemoEtAl2021}. Wu, Austrin, Pitassi, and
Liu prove SSE hardness for chordal Minimum Fill-In at every fixed
factor~\cite{APW}, but that result
does not cover trivially perfect completions, which form a more
restricted target class.
The adjusted-cost argument also transfers both constant total-depth
gaps to this completion objective, even on graphs of bounded maximum
degree: some factor $1+\epsilon_0>1$ is NP-hard, and every fixed
factor below $8/7$ is Unique Games-hard.

Section~\ref{sec:applications} derives the following consequences
from these constructions and from established reductions:
\Needspace{5\baselineskip}
\begin{itemize}
 \item \emph{Deletion and labeling objectives.} Vertex integrity
   minimizes the number of deleted vertices plus the largest remaining
   component. It inherits UGC hardness below $4/3$ and SSE hardness at every fixed
   factor. The SSE separator gap also gives this hardness for fracture
   number, starwidth, average and maximum general labels, and average
    hierarchical labels, using the comparisons
   in~\cite{GimaEtAlIntegrity} and the label lower bound
   of Blum and Storandt~\cite{BlumStorandt2023}.
   The SSE line-graph construction also gives edge-ranking hardness
   at every fixed factor on simple bipartite graphs, via
    $\operatorname{er}(B)=\td(\mathcal L(B))$~\cite{BodlaenderRankings}.
 \item \emph{Automata and hypergraph widths.} Gruber's established
   transformations~\cite{Gruber2008} carry the treedepth gaps to
   star height of regular languages represented by trim binary
   bideterministic partial automata. Replacing each vertex
   by a pair of adjacent copies carries the treewidth gaps to fractional, generalized,
   and ordinary hypertree width on rank-two hypergraphs, following
   the treewidth-hardness strategy of~\cite{Marx2010Fractional}.
 \item \emph{Junction-tree table size.} For binary variables, the
   total state space is hard to approximate within $2^{\chi n}$
   for every fixed $0<\chi<1/18$ under $\PP\ne\NP$, and for
   every fixed $0<\chi<1/6$ under UGC. These explicit exponents
   follow from the linear width gaps. Constant-factor hardness
   was proved by Bodlaender and Fomin~\cite{BodlaenderFomin2005}.
\end{itemize}
Established equivalent formulations of treedepth, including vertex
ranking and minimum symbolic elimination height, are covered by the
same parameter bounds rather than stated as separate results.

\subsection{Related work and attribution}

Bonnet's treewidth reduction~\cite{BonnetTreewidth} introduced the
truth-block and assignment-witness framework with three clique parts
used in Bonnet, Neuen, and Soko\l owski's treedepth reduction~\cite{BNS}.
Our NP-hardness proof inherits that encoding and their anchor graph.
We use their anchor theorem as a
black box, extending the anchor cliques only in the pairwise-incomparable
case and analyzing comparable anchors separately. Our changes are the
parity-witness stars and the partial-assignment analysis determining
the block sizes. Appendix~\ref{subsec:np-provenance} gives the detailed
comparison. The treewidth argument also uses Bonnet's central-bag
lower bound. Our replication lemma controls
the completeness overhead for the source gaps used here.

Khot and Regev~\cite{KhotRegev} supply the ordinary-UGC source for
the layered construction. B\'erczi, Chandrasekaran, Kir\'aly, Lee,
and Xu~\cite[Sections~3.5--3.6]{BCKLX2017} prove a weighted $k$-partite
Vertex Cover gap with ratio approaching $2(k-1)/k$. Our treewidth
transfer uses an equal-part unweighted gap, which we derive from
Khot--Regev by a multipartite product in
Section~\ref{sec:tw-transfer}. The
bounded-degree version uses the Gaussian density source of Austrin,
Khot, and Safra~\cite{AKS2011}. Our expander replacements use the
standard Margulis--Gabber--Galil family, as presented
in~\cite[Construction~8.1 and Theorem~8.2]{HLW2006}.

Wu, Austrin, Pitassi, and Liu~\cite{APW} use the expansion reductions
of Raghavendra, Steurer, and Tulsiani~\cite{RST} to prove SSE hardness
for treewidth, pathwidth, and related objectives at every fixed factor.
Their incidence-graph reduction also yields SSE treedepth hardness after a
separator-to-forest argument. We give that derivation in
Appendix~\ref{app:incidence-proof}.
For our
bounded-degree result, Section~\ref{sec:bounded-degree-hardness} also
derives the degree bound from the weighted source by the same
factor-completion tools used for NP and UGC hardness.

H\o gemo, Bergougnoux, Brandes, Paul, and Telle~\cite{HogemoEtAl2021}
study VPT-sum, relate it to trivially perfect completion, and prove
NP-hardness. The vertex-query interpretation described above gives
the same objective as expected search cost up to the factor
$|V(G)|$.
Szyfelbein~\cite{Szyfelbein2025GraphSearch} gives an
$O(\sqrt{\log n})$-approximation for the more general search
problem with nonuniform target weights and query costs. A subsequent
preprint also studies target-dependent costs~\cite{Szyfelbein2026GraphSearch}.
Our NP and UGC total-depth reductions have connected outputs. The UGC
graph $S_q(G)$ is connected whenever $G$ has an edge. The edgeless
source is always a YES instance and can be mapped to $K_1$. The NP
construction joins its three part cliques by source-edge paths. The
bounded-degree replacements retain these joining edges and use
connected sparse replacement graphs from Lemma~\ref{lem:factor-expanders}.
Thus the same
constant-factor hardness holds for average-case search with uniform
weights and unit query costs, also at bounded maximum degree.

The edge-query analogue is Dasgupta's similarity-based hierarchical
clustering objective~\cite{Dasgupta2016}. For a hierarchy $T$ with
leaves $V(G)$, its cost is
$\sum_{uv\in E(G)}w_{uv}\,|C_T(u,v)|$,
where $C_T(u,v)$ is the leaf set below the least common ancestor
of $u$ and $v$. For connected graphs, H\o gemo, Bergougnoux, Brandes,
Paul, and Telle~\cite[Theorem~6]{HogemoEtAl2021} identify its
unweighted optimum with EPT-sum, the minimum sum of leaf depths in an
edge-partition tree. The unweighted problem
is NP-hard by H\o gemo, Paul, and
Telle~\cite{HogemoPaulTelle2020}. Charikar and
Chatziafratis~\cite{CharikarChatziafratis2017} give an
$O(\sqrt{\log n})$-approximation and prove hardness within every
fixed constant factor under SSEH, including for unweighted graphs.
Szyfelbein and Dereniowski~\cite{SzyfelbeinDereniowski2026}
extend this SSE hardness to variants that stop at clusters inducing
trees or graphs of bounded diameter.
The vertex-query and edge-query formulations have different feasible
decompositions. The total-depth bounds proved here concern vertex
queries.

